\section{Rank-resolved receiver resources}
\label{sec:rankmemory92}
The distinction between classical and quantum receivers extends to an
attained hierarchy at each pair of retained-register dimensions.  This
section keeps the reset cut fixed and resolves its two sides separately.
It supplies a general variational characterization and an exactly solved
spectrum, rather than identifying all notions of bounded quantum memory.

\subsection{Two specified recording cuts}
Use the finite experiment of Definition~\ref{def:classes91}.
Fix integers $1\le k\le d_1$ and $1\le \ell\le d_2$.
\begin{definition}[Retained-register dimensions]\label{def:rankclass92}
The class $\mathfrak R_{k,\ell}$ consists of unit-reset strategies for
which, conditional on every recorded first history, the entire retained
old quantum system has dimension at most $k$ and the entire fresh
reference accompanying the second input has dimension at most $\ell$.
The old system and the complete fresh input--reference system are in a
tensor product at that cut.  All other quantum registers are discarded
before the second call; classical records may be retained without a
cardinality bound.  The initial first-call reference and temporary
workspaces before this cut have no dimension restriction.  The final
measurement on the two retained receivers is arbitrary.
Mixed preparations, recorded normal instruments, public randomization,
null histories, discarded terminal calls and finite-dimensional tester
closure have the conventions of Definition~\ref{def:classes91}.
Denote the optimal reward by $P_{k,\ell}$.
\end{definition}
In particular, this is not a bound on the largest workspace throughout
an implementation.  Enlarging the old output by keeping an instrument
environment would violate the stated resource.  The refinement in the
next lemma records that environment's Kraus index \emph{classically}
instead, without retaining it coherently.

\begin{lemma}[Dimension-preserving refinement]\label{lem:ranknormal92}
For the optimization over $\mathfrak R_{k,\ell}$, the normal form
\eqref{eq:resetgrams91}--\eqref{eq:resetblock91} is exact with
\begin{equation}\label{eq:rankconstraints92}
 \operatorname{rank}A_{yh}\le k,\qquad
 \operatorname{rank}b_{yh}\le\ell.
\end{equation}
Every implemented competitor can be refined into such branches while
preserving its score if the added classical indices are ignored in
its decisions.  Conversely every finite Gram collection satisfying
these ranks and the common-barycenter constraints is realizable in
$\mathfrak R_{k,\ell}$.
\end{lemma}
\begin{proof}
Purify the first input, retaining a reference whose relevant Schmidt
support has dimension at most $d_1$.  A completely positive branch of
the subsequent map from that support to $\C^k$ has a finite Kraus
representation $X\mapsto\sum_qV_{yhq}XV_{yhq}^*$, with
$V_{yhq}:\C^{r_0}\to\C^k$.  Refine $h$ to $(h,q)$, record $q$
classically, and use the original fresh preparation and final effects.
Summing over $q$ recovers the old channel and its decision statistics.
Thus $C_{yhq}=V_{yhq}C_0$ has rank at most $k$, and completeness
still gives $\sum_{h,q}C_{yhq}^*C_{yhq}=\rho$ for every $y$.
An original mixed first preparation is recovered by discarding its
purification before this old-output cut; this discarding is part of
that same completely positive branch, not an additional retained
quantum register.

A mixed fresh state on $I_2\otimes\C^\ell$ has a finite spectral
mixture of normalized pure states in that same space.  Sample and
record its mixture index, use the corresponding vector, and ignore
that index in the original decision.  Its coefficient map
$D:I_2\to\C^\ell$ has Gram rank at most $\ell$.  Multiplying the
old coefficient map by the square root of this mixing probability
preserves the common Gram sum.  This is a classical refinement, not a
purification into a larger fresh reference.  It is independent of the
old quantum output conditional on the refined record.  Countable
outcomes are handled by the trace tails of Lemma~\ref{lem:countable91}.
Public mixtures and tester limits preserve the resulting value bound.

Conversely choose $C_0^*C_0=\rho$ on its Schmidt support.  Factor each
$A_{yh}$ as $C_{yh}^*C_{yh}$ with
$C_{yh}:I_1\to\C^k$, and each $b_{yh}$ as $D_{yh}^*D_{yh}$ with
$D_{yh}:I_2\to\C^\ell$.  The maps $V_{yh}=C_{yh}C_0^+$ satisfy
$\sum_hV_{yh}^*V_{yh}=I$ on that support.  They give a recorded
trace-preserving instrument.  The normalized fresh vector $|D_{yh}\rangle$
is prepared independently, and polar isometries transfer the prescribed
leaf POVMs to $\C^k\otimes\C^\ell$.  Zero Gram branches are simply
zero maps.  The original unnormalized formulas define all null histories.
\end{proof}

\subsection{An attained hierarchy and its dual}
Write $\mathcal D_d^{\le q}$ for the compact set of density matrices of
rank at most $q$.  In the notation of \eqref{eq:leafvalue91}, put
\begin{align}
 g_y^\ell(a)&=\max_{b\in\mathcal D_{d_2}^{\le\ell}}\Phi_y(a,b),
 \label{eq:rankleaf92}\\
 c_y^{k,\ell}(\rho)&=\max\left\{\sum_h\lambda_hg_y^\ell(a_h):
 a_h\in\mathcal D_{d_1}^{\le k},\quad
 \lambda_h\ge0,\quad\sum_h\lambda_ha_h=\rho\right\}.
 \label{eq:rankroof92}
\end{align}
Weights sum to one by taking traces.  All density matrices have such a
representation, since rank-one atoms are allowed.

\begin{theorem}[Rank-resolved variational principle]\label{thm:rankroof92}
For every finite two-call experiment,
\begin{equation}\label{eq:rankprimaldual92}
\begin{split}
 P_{k,\ell}&=\max_{\rho\in\mathcal D_{d_1}}
                       \sum_yc_y^{k,\ell}(\rho)\\
 &=\min_{(H_y)}\lambda_{\max}\left(\sum_yH_y\right),\qquad
 \tr(H_ya)\ge g_y^\ell(a)\quad(a\in\mathcal D_{d_1}^{\le k}),
\end{split}
\end{equation}
where the $H_y$ are Hermitian.  Both extrema are attained.  At an
optimal barycenter $\rho$, at most $(\operatorname{rank}\rho)^2$
nonzero receiver outcomes after each first label suffice; a final
quantum receiver of dimension $k\ell$ suffices.
The values increase monotonically in either dimension, with
\[
 P_{1,d_2}=P_{\rm ca},\qquad P_{d_1,d_2}=P_{\rm reset}.
\]
The contact, spectral and quantitative defect criteria of
Theorem~\ref{thm:resetdual91} and Corollary~\ref{cor:contactdefect91}
hold with $g_y^\ell$ and the rank-restricted atoms.
\end{theorem}
\begin{proof}
The function $\Phi_y$ is continuous and homogeneous in its first
positive argument.  The rank-bounded sets are compact, so
$g_y^\ell$ is continuous and its graph on
$\mathcal D_{d_1}^{\le k}$ has compact convex hull.
The largest height over each $\rho$ is an attained, upper
semicontinuous concave envelope.  Lemma~\ref{lem:ranknormal92} gives
both operational directions.  The atom-elimination proof of
Theorem~\ref{thm:resetvariational91} uses only linear dependence of
atoms on the support of $\rho$, and does not change any atom.  It
therefore preserves the rank restriction and gives the same branch
bound.  Compactness supplies all fresh states and final POVMs.

For clarity, dual attainment does not require the rank-bounded atom
set to be convex.  The homogeneous hypograph of its concave envelope
is a closed convex cone, finite on every positive matrix and
nonnegative there.  At the common matrix $I/d_1$ with every height
$-1$, all these cones are strictly feasible.  The conic argument of
Theorem~\ref{thm:resetdual91} applies without alteration: its dual is
precisely majorization of the atom graph in \eqref{eq:rankprimaldual92}.
Bounded reward gives a finite value.  Slater duality yields no gap and
an attained Hermitian majorant.  Subtracting the primal value from the
dual upper gives the same three nonnegative defects, proving the
contact and quantitative assertions.  Inclusion of the rank-bounded
atom sets proves monotonicity.  The two endpoint identifications are
Theorems~\ref{thm:classicalroof91} and~\ref{thm:resetvariational91}.
\end{proof}
The dual inequalities are universal on their rank-bounded domains;
sampling is not a certificate.  In particular, strict improvement
between any two nested dimension pairs is equivalent to a feasible
finite strategy for the larger pair exceeding an attained dual upper
for the smaller pair.  This criterion covers arbitrary fixed pairs,
without asserting that every such pair exhibits a gap.

\subsection{Rank-sensitive swap bounds}
The rank in the next lemma is an upper bound on each of the two Gram
ranks; the active overlap may have a smaller rank.
\begin{lemma}[Swap bounds at finite rank]\label{lem:rankswap92}
Let $1\le r\le d$, $d\ge2$, and let $A,B\succeq0$ have ranks at most
$r$.  With $L=\sqrt A\otimes\sqrt B$,
\begin{align}
 \tr[-LSL]_+&\le\frac{r-1}{2r}\tr A\tr B,
 \label{eq:rankswapnegative92}\\
 \|L(I-dS)L\|_1&\le\left(d-\frac1r\right)\tr A\tr B.
 \label{eq:rankswapcenter92}
\end{align}
The bounds also hold if only the smaller of the two ranks is at most
$r$, and for rectangular factors with the same Grams.  Both bounds
are sharp: take $A=aP/r$, $B=bP/r$, where $P$ is any rank-$r$
projection and $a,b\ge0$.  For nonzero factors, equality in the first
bound with $r>1$ requires these forms with a common $P$; equality in
the second requires them when $(d,r)\ne(2,1)$.
\end{lemma}
\begin{proof}
The nonzero eigenvalues of $\sqrt A B\sqrt A$ number at most $r$.
Consequently, for $f=\|\sqrt A\sqrt B\|_1$,
\[
 \tr(AB)\ge f^2/r,\qquad f^2\le\tr A\tr B.
\]
Insert the first inequality into Lemma~\ref{lem:swapfilter90} to get
\eqref{eq:rankswapnegative92}.  The triangle step of
Lemma~\ref{lem:centeredswap91} gives
\[
 \|L(I-dS)L\|_1
 \le\tr A\tr B-\tr(AB)+(d-1)f^2
 \le\tr A\tr B+(d-1-r^{-1})f^2.
\]
The coefficient is nonnegative, proving \eqref{eq:rankswapcenter92}.
These arguments only use the smaller rank, not a common support
assumption.  The spectrum identities hold for singular factors by the
proved continuity argument, and polar isometries handle rectangular
maps.

For equality, when the coefficient of $f^2$ is positive, normalized
fidelity one gives $A/\tr A=B/\tr B$ without an invertibility
assumption, as in Lemma~\ref{lem:centeredswap91}.  Equality in the
rank-$r$ Cauchy--Schwarz inequality then forces exactly $r$ equal
nonzero eigenvalues, proving the common-projection forms.
For such forms the centered spectrum is $(1-d)ab/r^2$ with
multiplicity $r(r+1)/2$ and $(1+d)ab/r^2$ with multiplicity $r(r-1)/2$;
the negative swap has mass $ab(r-1)/(2r)$.  This proves sharpness.
The excluded centered case has zero coefficient and different equality
cases: for $d=2,r=1$, orthogonal pure factors also attain the bound.
For the first inequality at $r=1$, the negative mass is always zero.
\end{proof}

\subsection{The complete two-cut spectrum of the basis experiment}
The following formula resolves every pair of retained-register
thresholds on the same once-selected finite basis family.  No
additional design is introduced at a new threshold.
\begin{theorem}[Exact dimension spectra]\label{thm:rankspectrum92}
Fix $d\ge2$, $0\le t\le1$, and a weighted exact second-moment family
satisfying \eqref{eq:basismoments90}.  Under the once-drawn device rule,
let $P_{k,\ell}^{\rm b}$ use the priors of
Theorem~\ref{thm:operationalhierarchy90}, and let
$P_{k,\ell}^{\rm eq}$ use the equal hypothesis priors of
Theorem~\ref{thm:equalprior91}.  For $1\le k,\ell\le d$, set
$r=\min\{k,\ell\}$ and $L_0=2(d+1)-t^2$.  Then
\begin{align}
 P_{k,\ell}^{\rm b}
 &=\frac{d+1}{L_0}+\frac{t^2(r-1)}{2rL_0},
 \label{eq:rankbiased92}\\
 P_{k,\ell}^{\rm eq}
 &=\frac12+\frac{t^2(dr-1)}{2dr(d+1)}.
 \label{eq:rankequal92}
\end{align}
These are optima over the complete classes
$\mathfrak R_{k,\ell}$.  For $t>0$ they strictly increase at every
increase of the smaller threshold, and attain the unrestricted reset
values precisely at $k=\ell=d$.  The two-cut resource dependence is
therefore exactly through $\min\{k,\ell\}$ for these experiments.
\end{theorem}
\begin{proof}
Use the refined normal form of Lemma~\ref{lem:ranknormal92}; there is
no retained dilation environment outside the counted receivers.
Every branch has $\tr b_{yh}=1$ and smaller Gram rank at most $r$,
while $\sum_{yh}\tr A_{yh}=d$.
For the biased priors the equal-label payoff is
$-t^2 S/(dL_0)$ and all unequal-label payoffs are nonpositive.
Equation~\eqref{eq:rankswapnegative92} bounds the gain by
$t^2(r-1)/(2rL_0)$.
For equal priors write $\alpha=t^2/[2d^2(d+1)]$.
The $d-1$ opposite-sign payoff blocks cancel their denominator, so the
branch gain is at most
$\alpha\|(C_{yh}\otimes D_{yh})(I-dS)(C_{yh}\otimes D_{yh})^*\|_1$.
Equation~\eqref{eq:rankswapcenter92} and the Gram sum give
$\alpha d(d-1/r)$, which is \eqref{eq:rankequal92} minus $1/2$.
Classical mixtures and trace-norm closure preserve both upper bounds.

Here is a single finite construction attaining them.  Index the input
basis cyclically modulo $d$, and let $J_h:\C^r\to\C^d$ include the
$r$ consecutive coordinates starting at $h$, $0\le h<d$.
Put $P_h=J_hJ_h^*$; then $\sum_hP_h=rI$.
Choose the first coefficient map $C_0=I/\sqrt d$ and, after every
first device label, the receiver instrument and fresh coefficient maps
\begin{equation}\label{eq:cyclicrankprotocol92}
 V_h=J_h^*/\sqrt r,\qquad
 C_{yh}=J_h^*/\sqrt{dr},\qquad D_{yh}=J_h^*/\sqrt r.
\end{equation}
Thus $\sum_hV_h^*V_h=I$,
$A_{yh}=P_h/(dr)$, $b_{yh}=P_h/r$ and $\sum_hA_{yh}=I/d$.
The fresh vector has norm one and is independent of the old receiver.
Both retained quantum registers have dimension $r$; $h$ is classical.
The initial reference has dimension $d$, as explicitly allowed before
the recording cut in Definition~\ref{def:rankclass92}.

On these receivers the filtered swap is $S_r/(dr^2)$ and the centered
one is $(I_{r^2}-dS_r)/(dr^2)$.  For the biased prior, declare $C$
only on equal labels and the antisymmetric receiver projection.
For equal priors, declare $C$ on the antisymmetric projection on
equal labels and on the symmetric projection on unequal labels.
Complementary effects complete a POVM.  Their spectral signs attain
every branch bound above.  At $r=1$ the antisymmetric projection is
zero; the second readout reduces to the unequal-label rule.  At $r=d$
the instrument's redundant classical randomization may be omitted.
All constructions remain normalized at $t=0$ and on projective null
histories at $t=1$.
\end{proof}

\begin{corollary}[Adjacent gains and dimension certificates]
\label{cor:rankwitness92}
For $1\le r<d$, abbreviate the values in
Theorem~\ref{thm:rankspectrum92} by $P_r^{\rm b}$ and $P_r^{\rm eq}$.
Their adjacent gains are
\[
 P_{r+1}^{\rm b}-P_r^{\rm b}=\frac{t^2}{2L_0r(r+1)},\qquad
 P_{r+1}^{\rm eq}-P_r^{\rm eq}
       =\frac{t^2}{2d(d+1)r(r+1)}.
\]
For the equal-prior perturbations of Corollary~\ref{cor:gamestability91},
with its score error $\kappa$, every perturbed class value moves by at
most $\kappa$.  In particular an achieved ideal-reset score strictly
above $P_r^{\rm eq}+\kappa$ requires both counted receiver dimensions
to exceed $r$.  The adjacent gap survives if $2\kappa$ is smaller
than its displayed value.  A certified additional preparation defect
$\varepsilon$ changes the implemented score by at most $\varepsilon/2$
and may be added to these error budgets.
\end{corollary}
\begin{proof}
Subtract the consecutive exact formulas.  The hybrid and probability
law bounds are uniform in a fixed protocol, independently of its
receiver dimensions, so they hold for every rank-restricted supremum.
If either dimension were at most $r$, its nominal value would be at
most $P_r^{\rm eq}$ by Theorem~\ref{thm:rankspectrum92}.  This proves
the asserted certificate and the perturbation statements.
\end{proof}
These dimension certificates assume the specified reset architecture
and calibration bounds.  They are not device-independent conclusions
from an uncalibrated score.  The entire two-cut hierarchy still lies
inside the reset class; the unrestricted coherent-acquisition optimum
of Theorem~\ref{thm:equalprior91} is a different resource.

\subsection{Relation to constrained-memory discrimination}
Memory-constrained tester optimization and its relation to constrained
separability are established in \cite{OZNPQ89}, including the classical
adaptive boundary.  Theorem~\ref{thm:rankroof92} is a rank-envelope
characterization for the specified classical-output, fresh-preparation
cut, and Theorem~\ref{thm:rankspectrum92} solves its two-coordinate
values on one fixed finite family.  It does not replace the general
channel-memory formulations of that work.  The autonomous multi-time
hierarchy of \cite{ZB90} carries coherent information between successive
acquisitions and compiles arbitrary finite-time testers with a counter.
Here increasing the two receiver dimensions never removes the reset
cut.  Consequently saturation means the full reset optimum, not the
unrestricted strategy norm.  These different saturation statements
are compatible and cannot be interchanged.
